\documentclass[11pt,a4paper]{article}

\usepackage[utf8]{inputenc}
\usepackage[T1]{fontenc}
\usepackage{lmodern}
\usepackage[margin=2.7cm]{geometry}
\usepackage{amsmath,amssymb,mathtools}
\usepackage{booktabs}
\usepackage{enumitem}
\usepackage{xcolor}
\usepackage{hyperref}

\hypersetup{
    colorlinks=true,
    linkcolor=blue!50!black,
    citecolor=blue!50!black,
    urlcolor=blue!50!black
}

\setlength{\parskip}{0.65em}
\setlength{\parindent}{0pt}

\title{Merkblatt: XGBoost, Affine Smoothers, and Wald-AIPW Outcome Weights}
\author{Working note for the thesis}
\date{}

\begin{document}

\maketitle

\section{Purpose of this note}

This note summarizes why XGBoost is delicate in the outcome-weight framework of
Knaus (2024), especially for Wald-AIPW estimators. The aim is not to argue that
XGBoost is generally invalid as a nuisance learner in Double Machine Learning.
Rather, the point is more specific:

\begin{quote}
XGBoost may produce a reasonable Wald-AIPW point estimate, but the recovered
outcome weights may fail to be an exact finite-sample representation of that
estimate.
\end{quote}

This distinction is crucial. A DML point estimate can be computed with many
learners. However, Knaus-style outcome-weight diagnostics require stronger
algebraic conditions. In particular, the fitted outcome regressions must behave
like smoother operations on the outcome vector.

\section{Outcome weights and why smoothers matter}

The outcome-weight framework starts from estimators that can be written as
\begin{equation}
    \hat\tau = \omega'Y = \sum_{i=1}^N \omega_i Y_i .
    \label{eq:omega_rep}
\end{equation}

For pseudo-IV estimators, a sufficient route to such a representation is that the
pseudo-outcome can be written as a linear transformation of the observed outcome:
\begin{equation}
    \tilde Y = TY .
    \label{eq:pseudo_outcome}
\end{equation}
Then the estimator can be written as a weighted sum of observed outcomes.

For kappa estimators, this is usually straightforward. The pseudo-outcome has the
form
\begin{equation}
    \tilde Y_i = t_i Y_i,
\end{equation}
where \(t_i\) depends on \(Z_i\), \(D_i\), and the estimated propensity score, but
not on the outcome vector \(Y\). The corresponding matrix \(T\) is diagonal.

For Wald-AIPW and other DML estimators, the situation is different. The
pseudo-outcome includes fitted outcome regressions, such as
\begin{equation}
    \widehat{\mathbb E}[Y\mid Z=z,X].
\end{equation}
Therefore, the fitted outcome predictions themselves must be expressible as
smoother operations on \(Y\). Schematically, one needs
\begin{equation}
    \hat m = SY
    \quad\text{or more generally}\quad
    \hat m = a + SY .
    \label{eq:smoother}
\end{equation}

This is where XGBoost becomes delicate.

\section{Linear smoother versus affine smoother}

A \emph{linear smoother} has fitted values of the form
\begin{equation}
    \hat m = SY.
\end{equation}
A standard example is OLS with a fixed design matrix:
\begin{equation}
    \hat Y = X(X'X)^{-1}X'Y = HY,
\end{equation}
where \(H\) is the usual hat matrix.

An \emph{affine smoother} allows for an intercept-like component:
\begin{equation}
    \hat m = a + SY.
\end{equation}
The important point is not merely that the model contains a constant. The crucial
property for normalization is that the smoother treats constant shifts in the
outcome correctly. In particular, Knaus's affine-smoother condition can be
summarized as
\begin{equation}
    S\mathbf{1}_N = \mathbf{1}_N.
    \label{eq:affine_condition}
\end{equation}

In plain language, this means:

\begin{quote}
If every outcome is increased by the same constant \(c\), then every fitted
outcome prediction should also increase by \(c\).
\end{quote}

Equivalently, for shifted outcomes \(Y^c = Y+c\mathbf{1}_N\), a well-behaved
affine smoother satisfies
\begin{equation}
    \hat m(Y^c) = \hat m(Y) + c\mathbf{1}_N.
\end{equation}

This is why the condition matters for translation invariance. If the outcome
residual is
\begin{equation}
    Y_i - \hat m_i,
\end{equation}
then after shifting the outcome by \(c\), one wants
\begin{equation}
    (Y_i+c) - (\hat m_i+c) = Y_i-\hat m_i.
\end{equation}
The constant should cancel. If the fitted outcome prediction does not shift by
exactly \(c\), the constant does not cancel perfectly. Then the final estimator
can become sensitive to the outcome level.

\section{Scale-normalization and full normalization for Wald-AIPW}

For an exact outcome-weight vector \(\omega\), three sums are important:
\begin{align}
    C   &= \sum_{i=1}^N \omega_i, \\
    C_1 &= \sum_{i=1}^N \omega_i D_i, \\
    C_0 &= \sum_{i=1}^N \omega_i(1-D_i).
\end{align}

\subsection{Scale-normalization / translation invariance}

The estimator is translation invariant if
\begin{equation}
    \sum_{i=1}^N \omega_i = 0.
    \label{eq:sum_zero}
\end{equation}
This implies that additive outcome shifts do not change the estimate:
\begin{equation}
    \hat\tau(Y+c\mathbf{1}_N)-\hat\tau(Y)
    =
    c\sum_{i=1}^N\omega_i
    =
    0.
\end{equation}

For log outcomes, this is also connected to unit invariance. If wages are changed
from dollars to cents before taking logs, then
\begin{equation}
    \log(100Y_i) = \log(Y_i)+\log(100),
\end{equation}
which is just an additive shift of the log outcome.

For Wald-AIPW, scale-normalization requires the outcome smoother to satisfy the
affine condition. In the notation above, the key requirement is essentially
\begin{equation}
    S\mathbf{1}_N=\mathbf{1}_N.
\end{equation}
If this fails, even the weaker sum-to-zero property is not guaranteed.

\subsection{Full normalization}

Full normalization is stronger. It requires
\begin{equation}
    \sum_{i=1}^N \omega_iD_i = 1
    \quad\text{and}\quad
    \sum_{i=1}^N \omega_i(1-D_i) = -1.
    \label{eq:full_norm}
\end{equation}

If both conditions hold, the estimator behaves like a normalized treated-control
contrast:
\begin{equation}
    \hat\tau
    =
    \underbrace{\sum_{D_i=1}\omega_iY_i}_{\text{weighted treated mean}}
    +
    \underbrace{\sum_{D_i=0}\omega_iY_i}_{-\text{weighted untreated mean}}.
\end{equation}

Full normalization implies scale-normalization, because
\begin{equation}
    \sum_i\omega_i
    =
    \sum_i\omega_iD_i+\sum_i\omega_i(1-D_i)
    =
    1-1
    =
    0.
\end{equation}
But the reverse is not true. Wald-AIPW can be scale-normalized without being fully
normalized.

\section{Which conditions are needed for Wald-AIPW?}

For Wald-AIPW, there are two separate layers.

\subsection{Condition 3: affine outcome smoother}

The first important condition is the affine-smoother condition. In simplified
form:
\begin{equation}
    S\mathbf{1}_N=\mathbf{1}_N.
\end{equation}

This condition concerns the \emph{outcome regression}. It asks whether the learner
used for outcome prediction correctly reproduces constant shifts in the outcome.

If Condition 3 holds, the Wald-AIPW weights are expected to be scale-normalized:
\begin{equation}
    \sum_i\omega_i=0.
\end{equation}

If Condition 3 fails, even this basic property is not guaranteed. This is the
main XGBoost problem in the Knaus framework.

\subsection{Condition 5b: treatment predictions must use the same outcome smoothers}

Full normalization requires more than the affine outcome-smoother condition. For
Wald-AIPW, Knaus's additional condition can be summarized as follows.

Let \(S^z_1\) and \(S^z_0\) denote the smoother matrices used for the
instrument-specific outcome regressions. That is, these matrices generate fitted
outcome predictions such as
\begin{equation}
    \widehat{\mathbb E}[Y\mid Z=1,X],
    \qquad
    \widehat{\mathbb E}[Y\mid Z=0,X].
\end{equation}

For full normalization, the corresponding treatment predictions must be generated
by applying the same smoother matrices to the treatment vector \(D\):
\begin{equation}
    S^z_1D = \hat D^z_1,
    \qquad
    S^z_0D = \hat D^z_0.
    \label{eq:c5b}
\end{equation}

This is the condition that is easy to describe incorrectly. It is not that the
same smoother is used for ``treatment and instrument prediction''. The instrument
\(Z\) defines the groups. The condition is that, within the instrument-specific
regressions, the same smoother matrices used for the outcome predictions are also
used to construct the treatment predictions.

In plain language:

\begin{quote}
The treatment nuisance must be produced by applying the same instrument-specific
smoothers that were used for the outcome nuisance.
\end{quote}

This is usually not how standard flexible DML is implemented. Standard DML often
uses separate learners for outcome and treatment nuisance functions. Therefore,
Wald-AIPW may be scale-normalized but not fully normalized.

\section{Why random forests behave better than XGBoost}

A random forest prediction can often be represented as a weighted average of
observed outcomes:
\begin{equation}
    \hat m_i = \sum_{j=1}^N s_{ij}Y_j.
\end{equation}
If the forest weights satisfy
\begin{equation}
    \sum_{j=1}^N s_{ij}=1,
\end{equation}
then a constant shift in the outcome is reproduced exactly:
\begin{align}
    \hat m_i(Y+c\mathbf{1}_N)
    &=
    \sum_j s_{ij}(Y_j+c) \\
    &=
    \sum_j s_{ij}Y_j
    +
    c\sum_j s_{ij} \\
    &=
    \hat m_i(Y)+c.
\end{align}

Thus, forests naturally behave like affine smoothers when the forest weights are
row-normalized. This explains why ranger/random forest implementations often pass
the scale-normalization or translation-invariance diagnostic:
\begin{equation}
    \sum_i\omega_i=0.
\end{equation}

However, random forest Wald-AIPW may still fail full normalization if Condition
5b does not hold. Thus, the typical classification is:

\begin{center}
\begin{tabular}{lll}
\toprule
Learner & Condition 3 & Condition 5b \\
\midrule
Random forest / ranger & typically holds & usually fails in standard DML \\
XGBoost & can fail & usually fails in standard DML \\
\bottomrule
\end{tabular}
\end{center}

Hence:
\begin{align}
    \text{ranger Wald-AIPW}
    &\approx
    \text{scale-normalized but not fully normalized}, \\
    \text{XGBoost Wald-AIPW}
    &\approx
    \text{not guaranteed scale-normalized and not fully normalized}.
\end{align}

\section{What is \texttt{base\_score} in XGBoost?}

XGBoost does not start with an empty model. It starts with an initial prediction
for every observation. This initial prediction is called \texttt{base\_score}.

Schematically, the fitted XGBoost outcome regression has the form
\begin{equation}
    \hat m_i
    =
    \texttt{base\_score}
    +
    \nu f_1(X_i)
    +
    \nu f_2(X_i)
    +
    \cdots
    +
    \nu f_B(X_i),
    \label{eq:xgb_prediction}
\end{equation}
where \(f_b\) are boosted trees and \(\nu\) is the learning rate.

Thus, \texttt{base\_score} is the global starting value or intercept before any
tree corrections are added.

\subsection{Why automatic \texttt{base\_score} is delicate}

If \texttt{base\_score} is automatically estimated from the outcome, then the
initial intercept is itself a function of \(Y\):
\begin{equation}
    \texttt{base\_score} = g(Y).
\end{equation}

This creates an additional outcome-dependent component in the fitted values.
Unless the outcome-weight extraction explicitly accounts for this component, the
recovered smoother matrix may miss part of the dependence of \(\hat m\) on \(Y\).
Then the algebraic identity
\begin{equation}
    \hat\tau = \omega'Y
\end{equation}
may fail.

\subsection{Why fixed \texttt{base\_score = 0} is also not a guarantee}

One may try to avoid the hidden outcome-dependent intercept by fixing
\texttt{base\_score}, for example:
\begin{equation}
    \texttt{base\_score}=0.
\end{equation}

This makes the initial value known. But it does not guarantee the affine-smoother
condition. The reason is simple. If the outcome is shifted by \(c\),
\begin{equation}
    Y_i^c = Y_i+c,
\end{equation}
then the initial prediction remains
\begin{equation}
    \hat m_i^{(0)}=0.
\end{equation}
It does not automatically become \(c\). Therefore, the boosted trees must learn
the entire constant shift exactly.

With finite boosting rounds, shrinkage, regularization, and tree constraints,
there is no guarantee that XGBoost learns this constant shift perfectly. Hence,
one can have
\begin{equation}
    \hat m(Y+c\mathbf{1}_N)
    \neq
    \hat m(Y)+c\mathbf{1}_N.
\end{equation}

This is precisely a failure of the affine-smoother property.

\section{Why the other XGBoost hyperparameters matter}

In the implementation used for the learner comparison, the XGBoost learner was
configured with hyperparameters intended to make it closer to an affine smoother,
for example
\begin{verbatim}
alpha = 0
subsample = 1
max_delta_step = 0
base_score = 0
\end{verbatim}

The intuition behind these choices is as follows.

\begin{table}[h!]
\centering
\caption{Why the XGBoost hyperparameters matter for affine-smoother behavior}
\begin{tabular}{p{0.24\textwidth}p{0.66\textwidth}}
\toprule
Hyperparameter & Relevance for outcome-weight diagnostics \\
\midrule
\texttt{base\_score} &
Initial global prediction. If estimated automatically, it may introduce an
outcome-dependent intercept not captured by the recovered smoother. If fixed,
the trees must learn constant shifts exactly. \\
\addlinespace
\texttt{alpha = 0} &
Removes L1 regularization. L1 regularization can create thresholding and
non-smooth dependence of fitted values on \(Y\). \\
\addlinespace
\texttt{subsample = 1} &
Uses all observations for each tree. This removes one source of randomness and
makes the implicit smoother more stable. \\
\addlinespace
\texttt{max\_delta\_step = 0} &
Avoids imposing a cap on leaf updates. If updates are capped, constant shifts in
the outcome may not be absorbed correctly. \\
\addlinespace
\texttt{eta} / learning rate &
Shrinkage can prevent the model from absorbing constant shifts exactly unless
enough boosting rounds are used. \\
\addlinespace
\texttt{nrounds} &
With too few boosting rounds, the model may not learn the outcome level or its
constant shifts accurately. \\
\addlinespace
\texttt{lambda} &
L2 regularization shrinks leaf values and can prevent exact reproduction of
constant shifts. \\
\bottomrule
\end{tabular}
\end{table}

These choices may help, but they do not turn XGBoost into a guaranteed affine
smoother. Boosted trees remain an algorithmic learner whose fitted values are not
generally a row-normalized weighted average of observed outcomes.

\section{Why this can create strange Wald-AIPW results}

Wald-AIPW contains outcome residualization terms. Very schematically, it uses
objects of the form
\begin{equation}
    Y_i - \hat m_i.
\end{equation}

If the outcome smoother is affine, then after shifting \(Y_i\) by \(c\), the
prediction also shifts by \(c\), and the residual stays the same:
\begin{equation}
    (Y_i+c)-(\hat m_i+c)=Y_i-\hat m_i.
\end{equation}

If XGBoost does not reproduce the constant shift exactly, then
\begin{equation}
    (Y_i+c)-\hat m_i(Y+c)
    \neq
    Y_i-\hat m_i(Y).
\end{equation}
The constant level of the outcome leaks into the residualized term. Therefore,
the final estimator can become sensitive to arbitrary outcome shifts.

In outcome-weight language:
\begin{equation}
    \hat\tau(Y+c\mathbf{1}_N)-\hat\tau(Y)
    =
    c\sum_i\omega_i.
\end{equation}
If \(\sum_i\omega_i\neq0\), translation invariance fails.

This explains why XGBoost can produce strange values in Knaus's diagnostic
simulation. In that simulation, the artificial outcome is
\begin{equation}
    Y_i^* = 1 + D_i,
\end{equation}
so the true treatment effect is exactly one. A fully normalized estimator gives
\begin{equation}
    \omega'Y^*
    =
    \omega'\mathbf{1}_N + \omega'D
    =
    0+1
    =
    1.
\end{equation}
If the estimator is only scale-normalized but not fully normalized, then
\(\omega'\mathbf{1}_N=0\) but \(\omega'D\) need not equal one. If the estimator is
not even scale-normalized, then \(\omega'\mathbf{1}_N\neq0\), so the baseline
level \(1\) contaminates the estimated treatment effect.

This is why the XGBoost case is especially problematic in the outcome-weight
diagnostic framework.

\section{How this relates to the thesis implementation}

In the Vietnam application, the XGBoost Wald-AIPW point estimate is close to the
ranger estimate. This suggests that, in that low-dimensional near-random-assignment
design, the point estimate itself is not practically driven by the learner choice.

However, the algebraic identity check fails for XGBoost:
\begin{equation}
    \omega'Y \neq \hat\tau.
\end{equation}

This is the decisive diagnostic warning. It means that the recovered XGBoost
weights should not be interpreted as an exact outcome-weight representation of
the reported Wald-AIPW estimate. Therefore, diagnostics such as
\(\sum_i\omega_i\), effective sample size, negative weight shares, or covariate
balance under XGBoost weights must be interpreted cautiously.

The correct thesis interpretation is therefore:

\begin{quote}
XGBoost is included as a learner comparison for the DML point estimate. However,
because the algebraic identity \(\hat\tau=\omega'Y\) fails, the XGBoost weights are
not used as the preferred source of outcome-weight diagnostics. The main
diagnostic interpretation is based on learners for which the identity is verified,
such as linear outcome regression or ranger.
\end{quote}

\section{Short summary}

\begin{enumerate}[label=\arabic*.]
    \item Outcome weights require an exact representation
    \[
        \hat\tau=\omega'Y.
    \]

    \item For Wald-AIPW, this depends on how the outcome regressions are fitted.

    \item Condition 3 is the affine outcome-smoother condition:
    \[
        S\mathbf{1}_N=\mathbf{1}_N.
    \]
    It means that constant shifts in \(Y\) must be reproduced by the fitted
    outcome predictions.

    \item If Condition 3 holds, Wald-AIPW can be scale-normalized:
    \[
        \sum_i\omega_i=0.
    \]

    \item Full normalization is stronger and requires the treated and untreated
    outcome-weight sums to be \(+1\) and \(-1\).

    \item For Wald-AIPW full normalization, Condition 5b is also needed: the same
    instrument-specific outcome smoothers must be used to generate treatment
    predictions.

    \item Random forests often satisfy the affine-smoother condition because they
    predict by row-normalized weighted averages of observed outcomes.

    \item XGBoost can violate the affine-smoother condition because predictions
    start from \texttt{base\_score} and are then updated by boosted trees under
    shrinkage, regularization, and finite boosting rounds.

    \item Fixing \texttt{base\_score = 0} avoids a hidden automatically estimated
    intercept, but it does not guarantee affine behavior. The trees must still
    learn constant shifts exactly.

    \item Therefore, XGBoost may produce a reasonable DML point estimate but fail
    the exact Knaus outcome-weight identity:
    \[
        \omega'Y=\hat\tau.
    \]
\end{enumerate}

\section{One-sentence explanation}

\begin{quote}
The problem with XGBoost is not that it cannot be used to compute a Wald-AIPW
estimate; the problem is that its fitted outcome regression need not behave like
an affine smoother, so the recovered outcome weights may fail to represent the
estimate exactly and may not satisfy the normalization properties predicted by
Knaus's framework.
\end{quote}

\end{document}
